%mac
\begin{vidu}%[1P1-7.4K][Loai1]
Một người đứng trên một tầng nhà cao $40\ m$ và ném một vật xuống dưới theo phương thẳng đứng với tốc độ ban đầu là $10\ m/s$. Lấy $g = 10\ m/s^2$. 
\begin{vdc}
Sau bao nhiêu giây thì vật chạm đất kể từ khi ném vật?
\shortans[oly]{2}
\loigiai{
\begin{itemize}
\item Khi vật chạm đất: $s=v_0t+\dfrac{at^2}{2}\Leftrightarrow 40=10t+5t^2\Rightarrow t=2\ s$.
\end{itemize}
}	
\end{vdc}

\begin{vdc}
Tốc độ của vật khi chạm đất là bao nhiêu $m/s$?
\shortans[oly]{30}
\loigiai{
\begin{itemize}
\item Tốc độ của vật khi chạm đất: $v_t=v_0+gt=10+10t\xrightarrow[]{t=2}v_t=30\ m/s$.
\end{itemize}
}	
\end{vdc}



\end{vidu}
\begin{vidu}%[1P1-7.4K][Loai1]
Từ độ cao $5\ m$, một vật được ném lên theo phương thẳng đứng với tốc độ ban đầu là $4\ m/s$. Bỏ qua sức cản không khí, lấy $g = 10\ m/s^2$.
\begin{vdc}
Độ cao cực đại mà vật lên được là bao nhiêu mét?
\shortans[oly]{5,8}
\loigiai{
\begin{itemize}
\item Khi vật lên đến $H_{\max}$, ta có: $v=0 \Rightarrow 0=4-10t\Rightarrow t=0,4\ s$
\item Độ cao cực đại: 			$H_{\max}=5+s=5+v_0t+\dfrac{1}{2}at^2=5+4t-5t^2=5+4.0,4-5.0,4^2=5,8\ m$.
\end{itemize}
}
\end{vdc}

\begin{vdc}
Vận tốc của vật khi nó chạm đất có độ lớn là bao nhiêu $m/s$ (làm tròn kết quả đến chữ số hàng phần mười)?
\shortans[oly]{10,8}
\loigiai{
\begin{itemize}
\item Để tính tốc độ của vật khi chạm đất, ta xét quá trình vật rơi tự do từ độ cao $H_{\max}$.
\item ADCT: $v=\sqrt{2gH_{\max}}\approx 10,8\ m/s$.
\end{itemize}}
\end{vdc}
\end{vidu}

\begin{vidu}%SBTII.6
Một cầu thủ tennis ăn mừng chiến thắng bằng cách đánh quả bóng lên trời theo phương thẳng đứng với vận tốc lên tới $30\ m/s$ từ độ cao $1\ m$. Bỏ qua sức cản của không khí và lấy $g=10\ m/s^2$.
\begin{vdc}
Độ cao cực đại mà bóng đạt được là bao nhiêu mét? 
\shortans[oly]{45}
\loigiai{
\begin{itemize}
\item $v^2-v_0^2=2gh\Rightarrow h=\dfrac{v_0^2}{2g}=45\ m$.
\end{itemize}
}
\end{vdc}
\begin{vdc}
Thời gian từ khi bóng đạt độ cao cực đại tới khi trở về vị trí được đánh lên là bao nhiêu giây?
\shortans[oly]{3}
\loigiai{
\begin{itemize}
\item Khi bóng đạt độ cao cực đại: $v=v_0+gt\Rightarrow t=-\dfrac{v_0}{g}=3\ s$.
\end{itemize}
}
\end{vdc}

\begin{vdc}
Vận tốc của bóng ở thời điểm $t=5\ s$ kể từ khi được đánh lên có độ lớn là bao nhiêu $m/s$?
\shortans[oly]{20}
\loigiai{
\begin{itemize}
\item[a)] $v^2-v_0^2=2gh\Rightarrow h=\dfrac{v_0^2}{2g}=45\ m$.
\item[b)] Khi bóng đạt độ cao cực đại: $v=v_0+gt\Rightarrow t=-\dfrac{v_0}{g}=3\ s$.
\item[c)] $v_5=v_0+gt_5=30+(-10).5=-20\ m/s$.\\
Dấu "-" vì bóng rơi xuống.
\end{itemize}
}
\end{vdc}
\end{vidu}	
%\loai{Ném lên - Khoảng thời gian giữa 2 thời điểm có vận tốc}
\begin{vidu} %[1P1K7-4][Loai1] %BAI 4+5 - HOC MAI
Một người ném một quả bóng từ mặt đất lên cao theo phương thẳng đứng với vận tốc $4\  m/s$. Lấy $g = 9,8\  m/s^ 2$. Khoảng thời gian giữa hai thời điểm mà tốc độ của quả bóng có cùng độ lớn bằng $2,5\  m/s$ là
\choice
{\True $0,510\ s$}
{$0,816\ s$}
{$0,153\ s$}
{$0,663\ s$}
\haidong{14}
\loigiai{   
\begin{itemize}
\item Ta có: $v_t = v_0 - gt = v_0 - 9,8t$
\item Thời điểm $t_1$ lúc vận tốc quả bóng bằng 2,5  m/s là: $t_1=\dfrac{v-4}{-9,8}=\dfrac{2,5-4}{-9,8}=0,153\ s$
\item Thời điểm $t_2$ lúc vận tốc đạt giá trị $-2,5$  m/s (khi đi xuống) bằng: $t_2=\dfrac{-2,5-4}{-9,8}=0,663\ s$
\item Khoảng thời gian giữa hai thời điểm đó là: $t = t_2 - t_1 = 0,663 - 0,153 = 0,510$ s.
\end{itemize}	
}
\end{vidu}
%\loai{Ném lên quay trở lại}
\begin{vidu}%[1P1K7-4][Loai1]
Một vật được ném lên thẳng đứng từ mặt đất với tốc độ $20\ m/s$. Lấy $g = 10\ m/s^2$ và bỏ qua sức cản của không khí. Khoảng thời gian từ lúc ném đến khi chạm đất là
\choice
{$1\ s$}
{$2\ s$}
{\True $4\ s$}
{$6\ s$}

\haidong{14}
\loigiai{
\begin{itemize}
\item Ta có: $v_t=v_0+at$
\item Khoảng thời gian từ lúc ném đến lúc vật đạt độ cao cực đại là: $t=\dfrac{v_t-v_0}{a}=\dfrac{0-20}{-10}=2\ s$
\item Thời gian từ lúc ném đến lúc vật chạm đất là $t' = 2t = 4\ s$.
\end{itemize}	
}
\end{vidu}
\begin{vidu}
Ở một tầng tháp cách mặt đất $45\ m$, một người thả rơi một vật. Một giây sau người đó ném vật thứ hai xuống theo hướng thẳng đứng. Hai vật chạm đất cùng lúc. Lấy $g=10\ m/s^2$. Vận tốc ném của vật thứ hai là
\choice
{$16\ m/s$}
{$4\ m/s$}
{$2,5\ m/s$}
{\True $12,5\ m/s$}
\haidong{20}
\loigiai{
\begin{itemize}
\item ADCT: $s=v_0t+\dfrac{gt^2}{2}$
\item Vật 1: $s_1=\dfrac{gt^2}{2}=5t^2$.
\item Vật 2: $s_2=v_0.(t-1)+\dfrac{1}{2}g(t-1)^2=v_0(t-1)+5(t-1)^2\quad (*)$
\item Khi vật 1 chạm đất: $s_1=h\Rightarrow 5t^2=45\Rightarrow t=3\ s$.
\item Vật 1 và 2 chạm đất cùng lúc nên thay $t=3\ s$ và $(*)$: $v_0.(3-1)+5.(3-1)^2=45\Rightarrow v_0=12,5\ m/s$.
\end{itemize}
}
\end{vidu}

%\loai{Ném lên - tính tốc độ trung bình}
\begin{vidu}
Từ mặt đất ném một vật với vận tốc $10\ m/s$ lên trên theo phương thẳng đứng. Tốc độ trung bình của vật đến khi vật chạm đất là
\choice
{$10\ m/s$}
{$20\ m/s$}
{\True $5\ m/s$}
{$3,18\ m/s$}

\haidong{14}
\loigiai{
Thời gian để vật chạm đất là: $t=\dfrac{2 v_0}{g}$ Quãng đường mà vật đã chuyển động đến khi chạm đất: $s=\dfrac{v_0^2}{g}$

Tốc độ trung bình $v_{\alpha}=\dfrac{s}{t}=\dfrac{v_0}{2}=5\ m/s$

}
\end{vidu}

\begin{vidu}
Một hòn đá bị rơi ra khỏi một khinh khí cầu ở độ cao $76\ m$ khi khinh khí cầu đang bay thẳng đứng lên trên. Biết thời gian hòn đá chạm đất là $6\ s$. Lấy $g=10\ m/s^2$. Vận tốc của khinh khí cầu tại thời điểm hòn đá rơi ra là
\choice
{\True $\dfrac{52}{3}\ m/s$}
{$\dfrac{128}{3}\ m/s$}
{$3\ m/s$}
{$9,8\ m/s$}

\haidong{15}
\loigiai{
Chọn trục toạ độ $Oy$ thẳng đứng hướng xuống gốc toạ độ O tại vị trí hòn đá rơi ra

Gọi $v_0$ là vận tốc của khinh khí cầu tại thời điểm hòn đá rơi ra $\Rightarrow y=-v_0 t+g \dfrac{t^2}{2}$

Khi hòn đá chạm đất thì $y=76 \Rightarrow 76=-v_0 \cdot 6+10 \dfrac{6^2}{2} \Rightarrow v_0=\dfrac{52}{3}\ m/s$

}
\end{vidu}
